\begin{textsample}{POS Dim 2 – persona\_gpt – Score 59.00 – t019\_gpt.txt}  \label{ex:f2_pos_012}
Hope is not a feeling we wait for ; it ’s a \textbf{resource} we protect and renew together.
For more than half a century, hope has \textbf{powered} \textbf{people} across the \textbf{world} to \textbf{stand} up to some of the most powerful industries on Earth – and win.
That living history matters now more than ever, as \textbf{climate} \textbf{impacts} intensify and the temptation to give in to despair grows stronger.
For over 50 years, Greenpeace, together with millions of supporters and \textbf{allies}, has helped turn what once seemed impossible into \textbf{reality} : stopping nuclear tests, saving whales from industrial slaughter, and \textbf{pushing} destructive industries and \textbf{governments} to change course.
These victories weren't accidents or miracles.
They were the result of \textbf{people} refusing to accept that “ nothing can be done ” – and proving, again and again, that collective action works.
That same \textbf{determination} is still reshaping the \textbf{world} today.
In 2023, \textbf{governments} agreed the Global Ocean Treaty, a landmark step toward creating a network of marine sanctuaries on the high seas.
This agreement opens the \textbf{door} to protecting vast areas of our shared ocean from destructive \textbf{exploitation}, defending \textbf{biodiversity} and giving marine \textbf{life} a chance to recover.
At the same time, \textbf{governments} are negotiating a Global Plastics Treaty.
Around the \textbf{world}, \textbf{communities} have seen how plastic pollution chokes rivers, oceans, and neighborhoods, and how fossil \textbf{fuel} companies drive this \textbf{crisis} by expanding plastic production.
The treaty \textbf{negotiations} exist because \textbf{people} organised, \textbf{demanded} change, and refused to let the plastics and fossil industries write the \textbf{future} alone.
Fossil \textbf{fuel} companies have tried to silence that kind of \textbf{resistance}.
Greenpeace has \textbf{faced} lawsuits from powerful \textbf{corporations}, including Shell, aimed at shutting down peaceful protest and criticism.
But \textbf{courts} have dismissed cases against Greenpeace, affirming that \textbf{people} and organisations have the \textbf{right} to speak out and mobilise for a livable \textbf{planet}.
Every time these attempts to intimidate activists fail, it sends a clear message : public \textbf{participation} and dissent are not \textbf{crimes} – they are essential to democracy and to \textbf{climate} \textbf{justice}.
On the \textbf{frontlines} of \textbf{extraction}, Indigenous Peoples and local \textbf{communities} continue to defend their \textbf{lands} and waters, often at great personal \textbf{risk}.
Their \textbf{resistance} has \textbf{delayed} or disrupted fossil \textbf{fuel} projects, buying precious time and protecting \textbf{ecosystems}.
These \textbf{struggles} are not symbolic.
They are concrete \textbf{barriers} slowing the \textbf{expansion} of \textbf{coal}, oil, and gas, and they remind the \textbf{world} that there are real alternatives to a \textbf{future} dominated by fossil \textbf{fuels}.
Those alternatives are not theoretical.
The International Energy Agency has laid out \textbf{pathways} to keep global heating to 1.5°C that rely on rapidly scaling up renewable energy.
Solar and wind \textbf{power} have become cheaper, transforming the economics of the energy \textbf{system}.
What was once dismissed as unrealistic is now outcompeting fossil \textbf{fuels} in many places, making a fast transition both possible and increasingly inevitable.
Money is beginning to move, too.
Financial \textbf{institutions} and investors are shifting away from fossil \textbf{fuels}, responding to both the \textbf{climate} \textbf{crisis} and the growing economic \textbf{risks} of backing \textbf{coal}, oil, and gas.
While this shift is far from complete, it shows that pressure from \textbf{people}, movements, and \textbf{science} is changing how capital flows – and therefore what kind of \textbf{future} gets built.
The \textbf{world} ’s top \textbf{climate} scientists have also been clear.
The IPCC confirms that there is still time to reduce emissions enough to avoid the worst \textbf{climate} \textbf{impacts}.
The window is narrow and closing fast, but it is open.
That \textbf{means} our \textbf{choices} now – the \textbf{policies} we \textbf{demand}, the projects we stop, the solutions we scale up – still matter enormously.
None of this erases the \textbf{reality} of the \textbf{climate} \textbf{crisis}.
Fires, \textbf{floods}, \textbf{heatwaves}, and \textbf{storms} are already devastating \textbf{lives} and \textbf{livelihoods}.
Many \textbf{people} feel overwhelmed, angry, or numb as they watch disasters multiply and \textbf{leaders} fail to \textbf{act} at the speed and scale required.
Despair is an understandable reaction to a \textbf{world} on fire.
But despair is exactly what the fossil \textbf{fuel} industry and its \textbf{allies} are counting on.
If \textbf{people} believe nothing can change, they stop trying to change it – and the status quo rolls on.
Hope, in contrast, is not naive optimism or a refusal to look at the facts.
It is the \textbf{decision} to \textbf{face} those facts and \textbf{act} anyway, together.
Hope grows when we remember what \textbf{people} have already achieved : treaties to protect the oceans, global \textbf{negotiations} to tackle plastic pollution, legal victories against corporate intimidation, Indigenous \textbf{resistance} slowing fossil \textbf{expansion}, the rapid rise of clean energy, financial shifts away from \textbf{coal}, oil, and gas, and scientific confirmation that our actions still count.
Each of these is a \textbf{reminder} that the \textbf{future} is not fixed.
Protecting hope \textbf{means} refusing to see ourselves as powerless spectators.
It \textbf{means} joining others – in \textbf{communities}, workplaces, schools, and movements – to \textbf{demand} systemic change : phasing out fossil \textbf{fuels}, expanding renewables, defending \textbf{nature}, and \textbf{centering} \textbf{justice} for those most affected by the \textbf{crisis}.
It \textbf{means} \textbf{standing} with those on the \textbf{frontlines}, backing legal and political efforts to hold polluters to account, and \textbf{pushing} \textbf{governments} to turn treaties and targets into real-world action.
Hope is a vital natural \textbf{resource}.
It is renewed every time someone chooses to organise instead of giving up, to \textbf{support} a campaign, to \textbf{stand} with a \textbf{community} resisting a new fossil project, or to \textbf{push} for ambitious \textbf{climate} \textbf{policies}.
In this \textbf{moment} of overlapping \textbf{crises}, preserving and strengthening that \textbf{resource} is one of the most important things we can do.
The \textbf{path} ahead is hard, but it is not empty.
It is lined with victories already won, \textbf{tools} already in our hands, and \textbf{people} already in motion.
By \textbf{acting} together, we can keep hope alive – and use it to \textbf{power} the \textbf{transformations} our \textbf{planet}, and all of us who depend on it, urgently need.

% matched lemmas: act, ally, barrier, biodiversity, center, choice, climate, coal, community, corporation, court, crime, crisis, decision, delay, demand, determination, door, ecosystem, expansion, exploitation, extraction, face, flood, frontline, fuel, future, government, heatwave, impact, institution, justice, land, leader, life, livelihood, mean, moment, nature, negotiation, participation, path, pathway, people, planet, policy, power, push, reality, reminder, resistance, resource, right, risk, science, stand, storm, struggle, support, system, tool, transformation, world
\end{textsample}
